\subsection{由群在流形上的作用所定义的点分布场}

设 \((g,x)\mapsto\lambda(g,x)=gx\) 为 G 在 X 上的一个 \(C^{r}\) 类左作用律。设 \(s\leqslant r\) 且 \(t\in U_{s}(G)\)。对所有 \(x\in X\)，有 \(t*\varepsilon_{x}\in T_{x}^{(s)}(X)\)。映射 \(x\mapsto t*\varepsilon_{x}\) 称为由 t 及 G 在 X 上的作用所定义的点分布场，有时记作 \(D_{t}^{\lambda}\) 或简记作 \(D_{t}\)。设 \(\Omega\) 为 X 的开子集，F 为分离多范空间。若 \(f\colon\Omega\to F\) 是 \(C^{r}\) 类的且 \(s\leqslant r\)，则函数 \(t^{\vee}\ast f\)（\(\Omega\) 上的）也记作 \(D_{t}.f\)。于是

\[
(D _ {t} f) (x) = \langle t * \varepsilon_ {x}, f \rangle .\tag{16}
\]

若 \(s < \infty\)，则由第 4 节，\(D_t f \in \mathcal{C}^{r-s}(\Omega, F)\)。这样，\(f \mapsto D_t f\) 是 \(\mathcal{C}^r(\Omega, F)\) 到 \(\mathcal{C}^{r-s}(\Omega, F)\) 的映射（由于记号的滥用，常记作 \(D_t\)）。

若 \(t \in \mathrm{U}_s(\mathrm{G})\)、\(t' \in \mathrm{U}_{s'}(\mathrm{G})\) 且 \(s + s' \leqslant r\)，则由第 4 节的命题 18，

\[
\mathrm{D} _ {t * t ^ {\prime}} f = \mathrm{D} _ {t ^ {\prime}} (\mathrm{D} _ {t} f)\tag{17}
\]

因而，利用上面指出的记号滥用，有

\[
\mathrm{D} _ {t * t ^ {\prime}} = \mathrm{D} _ {t ^ {\prime}} \circ \mathrm{D} _ {t}.\tag{18}
\]

设 G 与 X 都是有限维的。映射 \((t, x) \mapsto t \otimes \varepsilon_{x}\)（从 \(\mathrm{T}^{(\mathrm{s})}(\mathrm{G}) \times \mathrm{X}\) 到向量丛 \(\mathrm{T}^{(\mathrm{s})}(\mathrm{G} \times \mathrm{X})\)，cf.~Differentiable and Analytic Manifolds, R, 13.2.5）是 \(C^{r-s}\) 类的。因而 (Differentiable and Analytic Manifolds, R, 13.2.5)，映射 \((t, x) \mapsto t * \varepsilon_{x}\)（从 \(\mathrm{T}^{(\mathrm{s})}(\mathrm{G}) \times \mathrm{X}\) 到向量丛 \(\mathrm{T}^{(\mathrm{s})}(\mathrm{X})\)）是 \(C^{r-s}\) 类的。特别地，\(D_{t}\) 是 Differentiable and Analytic Manifolds, R, 14.1.6 意义下阶数 \(\leqslant s\) 且类为 \(C^{r-s}\) 的微分算子。由公式 (16)，函数 \(D_{t}f\) 就是 f 在这个微分算子下的像 (Differentiable and Analytic Manifolds, R, 14.1.4)。

我们现在不再假设 G 与 X 都是有限维的。设 \(\psi\) 为流形 X 的自同构，\(\Delta\) 为 X 上的一个点分布场。按照一般的定义，\(\Delta\) 在 \(\psi\) 下的变换是 X 上的这样一个点分布场：它在 \(\psi(x)\) 的值是 \(\psi_{*}(\Delta(x))\)；我们把这个映射记作 \(\psi(\Delta)\)。若 \(g \in G\) 且 \(\tau(g)\) 表示 X 的自同构 \(x \mapsto gx\)，则 \(\Delta\) 在 \(\tau(g)\) 下的变换也称为 \(\Delta\) 在 g 下的变换。

命题 21. 设 \(\psi\) 为与 G 的作用交换的 X 的自同构。则 \(D_{t}\) 在 \(\psi\) 下不变。

对所有 \(x \in X\)，

\[
\begin{array}{r l} (\psi (\mathrm{D} _ {t})) (\psi (x)) & = \psi_ {*} (\mathrm{D} _ {t} (x)) = \psi_ {*} (t * \varepsilon_ {x}) \\ & = t * \psi_ {*} (\varepsilon_ {x}) \quad \text {(Proposition 15)} \\ & = t * \varepsilon_ {\psi (x)} = \mathrm{D} _ {t} (\psi (x)). \end{array}
\]

命题 22. 若 \(g \in G\)，则 \(D_t\) 在 \(g\) 下的变换是 \(D_{\varepsilon_g * t * \varepsilon_{g^{-1}}}\)。

这个变换在 gx 处的值是

\[
\begin{array}{r l r} \tau (g) _ {*} (\mathrm{D} _ {t} (x)) & = \tau (g) _ {*} (t * \varepsilon_ {x}) \\ & = \varepsilon_ {g} * (t * \varepsilon_ {x}) & \text {(Proposition 14 (i))} \\ & = (\varepsilon_ {g} * t * \varepsilon_ {g^{-1}}) * \varepsilon_ {g x} & \text {(Propositions 1 and 2)} \\ & = \mathrm{D} _ {\varepsilon_ {g} * t * \varepsilon_ {g^{-1}}} (g x). \end{array}
\]

设 \((x,g)\mapsto \mu (x,g) = xg\) 为 G 在 X 上的一个 \(\mathbf{C}^r\) 类右作用律。设 \(s\leqslant r\) 且 \(t\in \mathrm{U}_{s}(\mathrm{G})\)。对所有 \(x\in \mathbf{X}\)，有 \(\varepsilon_{x}*t\in \mathrm{T}_{x}^{(s)}(\mathbf{X})\)。映射 \(x\mapsto \varepsilon_x*t\) 称为由 \(t\) 及 G 在 X 上的作用所定义的分布场，有时记作 \(\mathrm{D}_{t}^{\mu}\) 或简记作 \(\mathrm{D}_t\)。设 \(\Omega\) 为 X 的开子集。若 \(f:\Omega \to F\) 是 \(\mathbf{C}^r\) 类的，则函数 \(f*t^{\vee}\) 记作 \(\mathrm{D}_t f\)。于是

\[
\left(\mathrm{D} _ {t} f\right) (x) = \langle \varepsilon_ {x} * t, f \rangle\tag{19}
\]

且在明显的记号下，有

\begin{enumerate}
\def\labelenumi{(\arabic{enumi})}
\setcounter{enumi}{19}
\tightlist
\item
\end{enumerate}

\[
D _ {t * t ^ {\prime}} f = \mathrm{D} _ {t} (\mathrm{D} _ {t ^ {\prime}} f)\tag{20}
\]

\[
\mathrm{D} _ {t * t ^ {\prime}} = \mathrm{D} _ {t} \circ \mathrm{D} _ {t ^ {\prime}}\tag{21}.
\]

命题 21 仍然成立。设 \(g \in G\)。\(D_t\) 在 \(g\) 下的变换（即在自同构 \(x \mapsto xg\)（\(X\) 的自同构）下的变换）是 \(D_{\varepsilon_{g^{-1}} * t * \varepsilon_g}\)。

